% ============================================================
\chapter{Literature Review}
% ============================================================

\section{Review of Literature}

\para{Traditional Network Management}

Network management has long relied on standard protocols rather than on
natural language. The Simple Network Management Protocol (SNMP) lets a
management station read and change variables on routers, switches and other
devices (Case et al., 1990), and its later architecture added security and
extensibility to the same model (Harrington et al., 2002). NETCONF replaced
much of SNMP's configuration role with a structured, transaction-based
protocol for installing and changing device configuration (Enns et al.,
2011). These protocols made automation possible, but they expose low-level
data models: an operator, or a script written by one, must still know
exactly which variable or configuration element to read or change.

Intent-Based Networking (IBN) is the research community's response to this
gap. RFC~9315 defines intent as a statement of the outcome an operator
wants, rather than the steps to reach it, and describes systems that
translate such intent into configuration and then check that it is being met
(Clemm et al., 2022). The same document notes that intent-based networking is
a natural candidate for artificial intelligence techniques. LLMs offer a
practical way to express intent in everyday language, which is the approach
NetMind takes.

\para{Language Models That Use Tools}

NetMind depends on an LLM being able to act, not only to answer. Yao et al.\
(2023) showed with ReAct that interleaving reasoning steps with actions such
as calls to external tools makes language models more accurate and their
behaviour easier to inspect, and Schick et al.\ (2023) showed with Toolformer
that a model can learn when and how to call external interfaces on its own.
Modern LLM APIs expose this ability as \emph{function calling} or \emph{tool
calling}, which NetMind uses: the model chooses a tool and its arguments,
and NetMind's own code decides whether and how to run it.

\para{LLMs for Network Management}

Research on applying large language models (LLMs) to network management has
grown rapidly since 2023. Hong, Van Tu, and Hong (2025) provide a
comprehensive survey of this field, classifying LLM-based approaches into
five task areas: network design, configuration, fault management, security,
and orchestration. Their analysis shows that most existing studies address
only one stage of the management workflow in isolation, rely on simulated or
historical data rather than live infrastructure, and still suffer from
hallucination and limited reliability when scaled to complex, heterogeneous
environments.

At the systems level, Wang et al.\ (2025) describe Confucius, a multi-agent
LLM framework developed at Meta that models network management workflows as
directed acyclic graphs and coordinates several specialised agents, each
responsible for a distinct sub-task, to carry out complex operator intents
while validating actions before execution. This illustrates one architectural
direction for LLM-based network management: decomposing a task among
cooperating agents rather than relying on a single general-purpose model.

Closer to the security-focused component of NetMind, Jiang, Gember-Jacobson,
and Feamster (2024) propose CAIP, a framework that supplies LLMs with
carefully mined network-specific context to detect router misconfigurations,
reporting a full detection rate on synthetic faults and the discovery of
previously unnoticed issues in real configurations. Lorusso, Maci, and Coscia
(2025) propose GOLLUM, a conversational agent that guides firewall
configuration using a retrieval-augmented LLM to compensate for the limited
domain-specific knowledge of general-purpose models. Both studies confirm that
LLMs can meaningfully assist with network security tasks when given the right
supporting context, but each addresses a single, narrow task rather than an
integrated management platform.

\para{AI on Both Sides of Security}

The same abilities are changing cybersecurity in both directions. A
systematic review of 185 studies on LLMs for cybersecurity identifies
LLM-based autonomous agents as a significant emerging trend, moving from
single tasks to orchestrating complex, multi-step security workflows (Xu et
al., 2024). On the attacking side, Fang et al.\ (2024) showed that a single
LLM agent given tools and a vulnerability description could autonomously
exploit 87\% of a benchmark of 15 real-world one-day vulnerabilities, where
open-source vulnerability scanners exploited none. The incidents described
in Chapter~1, the agent-driven intrusion at Hugging Face and the AI-assisted
PHP2Shell exploits, show that this is no longer only a laboratory result.
Defenders therefore need tools that work at a comparable speed, but because
a wrong change to a network can itself cause an outage, those tools must
keep a human in control of every change. This balance is the design goal of
NetMind.

\section{Brief Description of Existing Systems}

Beyond academic studies, network management is supported by established
monitoring tools and, more recently, by LLM-based assistants offered as
commercial products and independent open-source projects.

Nagios (Nagios Enterprises, n.d.) and Zabbix (Zabbix LLC, n.d.) are widely
used open-source monitoring systems. They poll devices and services, often
over SNMP, show their state on dashboards, and raise alerts when a check
fails or a threshold is crossed. They are mature and scale to large networks,
but each check, threshold and alert must be configured by an expert, they
have no natural-language interface, and acting on an alert still means
returning to the command line.

Selector AI Copilot (Selector AI, n.d.) is a commercial platform that
grounds an LLM in real network telemetry, allowing operators to ask natural-language questions about
their infrastructure and, through an agent framework, translate the answers
into automated operational workflows.

JACE (ntunes, n.d.) is an open-source autonomous agent that connects to
real Junos network devices over SSH and NETCONF, performs scheduled health checks, detects
anomalies, and answers natural-language questions about the network, while
retaining memory of device state and past incidents across sessions.

The two LLM-based systems represent, respectively, an enterprise-grade
commercial product and an independently developed open-source project, both
operating exclusively against real production networks.

\section{Advantages and Disadvantages of Existing Systems}

The systems reviewed above demonstrate clear strengths. Traditional
monitoring tools such as Nagios and Zabbix are proven, free to use and
reliable at scale, and they detect problems without anyone having to ask.
Selector AI benefits
from mature integration with real telemetry sources and an established agent
framework suited to production environments. JACE benefits from persistent
memory across sessions, allowing it to correlate incidents and remember
device baselines over time, and from directly automating monitoring without
requiring an operator to ask. CAIP and GOLLUM demonstrate that narrowly scoped
LLM systems can achieve very high accuracy within their specific task, since
focusing on one problem allows deeper, more reliable context engineering.

These strengths come with trade-offs. Selector AI is designed for large, real
production networks and offers no safe environment in which students or new
operators can experiment without risk to a live network. JACE, and most of
the studies surveyed by Hong et al.\ (2025), operate only against real
devices, so an operational mistake by the LLM or the operator directly affects
real infrastructure; none of them provide an isolated simulated environment
for safe testing. CAIP and GOLLUM are limited to a single task,
misconfiguration detection or firewall configuration respectively, and do not
extend to device control, monitoring, or other management functions, so they
would need to be combined with other tools to form a complete management
platform.

There is no system among those reviewed that combines a safe testing
environment with broad, integrated functionality while remaining open source.

\section{Problems with Existing Systems}

Drawing on the systems above, four recurring problems can be identified.
The first three concern current LLM-based network management tools. First, most systems that act on
real infrastructure provide no safe, isolated environment in which an LLM's
actions can be tested before being trusted, which is a significant barrier for
an educational or early-stage project where mistakes must be recoverable.
Second, systems that do provide broad functionality, such as Confucius,
achieve this through complex multi-agent orchestration, which increases
engineering effort and is difficult to replicate at a smaller scale. Third,
task-specific systems such as CAIP and GOLLUM achieve strong results within
their domain but do not generalise to device control, network visibility, or
camera and endpoint security, leaving a gap between narrow, highly accurate
tools and broad, general-purpose but less specialised ones. Fourth,
traditional monitoring tools report problems well but cannot be asked
questions in everyday language and do not help the operator act on what
they find, which keeps the cognitive load described in Chapter~1 on the
operator.

{\small
\begin{xltabular}{\textwidth}{|L{3.3cm}|L{2.7cm}|Y|L{2.1cm}|L{1.7cm}|}
\hline
\head \textbf{System (type)} & \textbf{Safe testing environment} & \textbf{Task scope} & \textbf{Natural language} & \textbf{Open source} \\ \hline
\endhead
Nagios / Zabbix (traditional monitoring) & No: monitors live networks & Broad: polling, dashboards and threshold alerts & No & Yes \\ \hline
Selector AI Copilot (commercial AIOps platform) & No: live production networks only & Broad: observability and agent workflows & Yes & No: commercial \\ \hline
JACE (open-source agent for Junos) & No: real devices over SSH / NETCONF & Broad: health checks, anomalies, Q\&A & Yes & Yes \\ \hline
CAIP (research prototype) & N/A: research only & Narrow: router misconfiguration detection & No: LLM used internally & Partial \\ \hline
GOLLUM (research prototype) & N/A: research only & Narrow: firewall configuration & Yes & Partial \\ \hline
Confucius (multi-agent LLM framework) & No & Broad but complex: multi-agent orchestration & Yes (operator intents) & No \\ \hline
\textbf{NetMind (this project)} & \textbf{Yes: simulated lab and read-only real mode} & \textbf{Broad: twenty-one tools for devices, firewalls, cameras, diagnostics and discovery} & \textbf{Yes: Arabic and English} & \textbf{Yes: on GitHub} \\ \hline
\caption{Comparison of NetMind with existing systems}\label{tab:comparison}\\
\end{xltabular}}

\section{Proposed System}

NetMind addresses these problems directly. It pairs a fully simulated
Containerlab network, in which every device is a real, independently
controllable Docker container, with a strictly read-only real-network
discovery mode, so that destructive or exploratory actions by the LLM are
confined to a disposable simulated environment while the real network can
only ever be observed. The boundary is enforced in code: the real-network
backend simply has no capability to change a device, whatever the model
asks. This removes the safety concern identified in existing
real-network-only tools such as Selector AI and JACE. Rather than a
multi-agent architecture, NetMind uses a single language model with a
defined set of callable tools spanning device control, network visibility,
connectivity checks, firewall auditing, and camera security, which keeps the
system simpler to build and reason about than a framework like Confucius
while still covering more ground than single-task systems such as CAIP or
GOLLUM. Each conversation is paired with a live dashboard and an alert
history that reflect the current state of the network, so that every query
is answered from up-to-date system state rather than from potentially stale
conversational memory. Unlike traditional monitoring tools, the same
conversation that reports a problem can also, after the operator confirms,
act on it.

\section{Alternative Solutions Considered}

\para{Network Simulation Platform}

Several network simulation platforms were evaluated before choosing
Containerlab for NetMind's simulated network. Cisco Packet Tracer is widely
used for teaching, but its nodes are not real operating systems and cannot be
controlled programmatically from outside the application. GNS3 and EVE-NG run
real device images, but typically inside full virtual machines, which is
heavier on resources and oriented toward a graphical, VM-centric workflow.
Mininet is lightweight and Python-based, but its network model is built
around virtual switches for software-defined-networking research rather than
independently controllable general-purpose containers. Kathara is
Docker-based, similar in spirit to Containerlab, but has a smaller ecosystem
and less adoption in real network-engineering contexts. Containerlab was
chosen because each node it creates is a real, independently addressable
Docker container: NetMind's control layer can use the standard Docker SDK to
start, stop, and execute commands inside any device exactly as on a real
Linux host, including running iptables on the firewalls and a small HTTP
service on the cameras, a level of direct, programmatic, container-native
control that none of the other evaluated tools offered.

\para{Large Language Model Provider}

Several ways of providing the language model behind NetMind were also
considered. The OpenAI and Anthropic APIs offer strong general-purpose models
with mature tool-calling support, but at higher per-token latency and cost
than were necessary for this project. A locally self-hosted model, for
example through Ollama, would avoid ongoing API costs, but requires capable
local hardware to run reliably enough for consistent tool-calling and
bilingual Arabic/English understanding, which was not practical in the
project's development environment. Groq was chosen because it exposes an
OpenAI-compatible tool-calling interface while running open-weight models
such as \code{openai/gpt-oss-120b} on inference hardware built specifically
for low response latency. Since NetMind's agent loop can make several
sequential tool calls to satisfy a single user request, per-call latency
compounds directly into the user's perceived response time, making Groq's
inference speed a directly relevant advantage rather than an incidental one.
The model is reached through a small provider interface (Section~5.2.3), so
this choice can be revisited without changing the rest of the system.
